\section{Introduction}
\label{sec:intro}

A user searching for the monthly rent of a specific apartment, the latest release
of a software package, or the current price of a product expects the search result
to reflect what the web page says \emph{today}. In practice, search engine indices
are built from crawls that may be days or weeks old~\citep{cho2003estimating}, and
the underlying pages continue to change between crawls. The result is a systematic
mismatch: the answer embedded in a search snippet or a retrieval-augmented generation
(RAG) response may be factually correct as of the last crawl but wrong at query time.

This \emph{temporal staleness} problem is qualitatively familiar---practitioners
routinely observe that search results lag behind price movements, software releases,
and real-estate inventory---but it resists quantitative study for a simple reason:
there is no standard benchmark. Without ground-truth change labels, it is impossible
to measure recall at fresh-versus-stale granularity, to compare retrieval systems on
time-sensitive queries, or to train models that detect whether a cached result is
still accurate.

We introduce \textbf{FreshState}, a benchmark construction pipeline designed to
fill this gap. FreshState operates on two levels. At the \emph{data collection}
level, it deploys a prospective monitor that fetches candidate URLs daily, extracts
answer-bearing values (prices, version strings) via domain-specific extractors, and
records change events when those values shift between consecutive snapshots. At the
\emph{evaluation} level, it pairs each change event with a stale search snippet and
measures whether language models echo the outdated value---a controlled
\emph{snippet-swap} experiment that quantifies the downstream impact of index lag.

\paragraph{Contributions.} We make the following contributions:

\begin{enumerate}
  \item \textbf{FreshState pipeline.} A modular, domain-extensible pipeline for
        constructing temporally-grounded web staleness benchmarks, covering candidate
        collection, baseline snapshot, daily monitoring, and change-event extraction
        (\Cref{sec:method}).

  \item \textbf{Domain-specific extractors.} A library of extractors for prices,
        availability signals, version strings, and GitHub release tags, with a shared
        confidence-scoring interface that enables automated quality filtering.

  \item \textbf{Large-scale prospective deployment.} Results from a 14-day deployment
        monitoring 1,628 URLs across two high-churn domains (rental listings and
        open-source software), producing 474 change-event seeds
        (\Cref{sec:experiments}).

  \item \textbf{Snippet-swap evaluation.} A controlled experiment showing that
        language models echo stale values at high rates when given outdated search
        context, demonstrating that index lag directly degrades answer quality in
        RAG pipelines (\Cref{sec:snippet_swap}).
\end{enumerate}

\paragraph{Findings.} Software releases change at a rate of ${\sim}17\%$ of tracked
repositories per day, producing 146 change events in the first 5 days alone.
Apartment listings exhibit lower value-change rates (${\sim}1\%$/day) but extremely
high expiration rates: 23\% of URLs return HTTP~410 (Gone) within 24 hours of
baseline collection, meaning any search engine that indexed them even one day earlier
would serve a dead link. In the snippet-swap experiment, stale context causes models
to reproduce the outdated value in 100\% of cases, confirming that
temporal staleness is not merely an index-freshness problem but a direct source of
factual error in language model outputs.

The rest of this paper is organized as follows. \Cref{sec:related} surveys temporal
IR, web archiving, and fact-verification work. \Cref{sec:method} formalizes the
problem and describes the FreshState pipeline. \Cref{sec:experiments} presents
dataset construction experiments and change-rate statistics. \Cref{sec:snippet_swap}
describes the snippet-swap evaluation. \Cref{sec:conclusion} discusses limitations
and future directions.
